% Ce chapitre repose sur :
%   DOC: s06, s07, s08, s09, s10
%   OBS: REL-DEM.I01, REL-DEM.I02, REL-DEM.P01, REL-DEM.P02, REL-DEM.P03, REL-DEM.P04, REL-DEM.P05,
%   OBS: REL-DEM.M01, REL-DEM.M02, REL-DEM.M03, REL-DEM.M04, REL-IPP.O01, REL-IPP.O02, REL-IPP.O03,
%   OBS: REL-IPP.O04, REL-IPP.C01, REL-IPP.C02, REL-IPP.C03, REL-IPP.C04, REL-IPP.C05, REL-IPP.C06,
%   OBS: REL-IPP.C07, REL-IPP.C08, REL-IPP.C09, REL-IPP.C10, REL-IPP.C11, REL-IPP.C12, REL-IPP.C13,
%   OBS: REL-IPP.F01, REL-IPP.F02, REL-IPP.F03, REL-IPP.F04, REL-IPP.R01, REL-IPP.R02, REL-IPP.R03,
%   OBS: REL-IPP.R04, REL-IPP.R05, REL-IPP.R06, REL-IPP.R07, REL-IPP.R08, REL-IPP.R09, REL-IPP.R10,
%   OBS: REL-IPP.R11, REL-IPP.R12, REL-PAT.T01, REL-PAT.T02, REL-PAT.T03, REL-PAT.D01, REL-PAT.D02,
%   OBS: REL-PAT.D03, REL-PAT.D04, REL-PAT.D05, REL-PAT.D06, REL-PAT.D07, REL-PAT.D08, REL-PAT.D09,
%   OBS: REL-PAT.D10, REL-PAT.D11, REL-PAT.D12, REL-PAT.D13, REL-PAT.D14, REL-PAT.D15, REL-PAT.D16,
%   OBS: REL-PAT.A01, REL-PAT.A02, REL-PAT.A03, REL-PAT.A04, REL-PAT.A05, REL-PAT.A06, REL-PAT.H01,
%   OBS: REL-PAT.H02, REL-PAT.H03, REL-PAT.H04, REL-PAT.H05, REL-PAT.H06, REL-PAT.H07, REL-PAT.H08,
%   OBS: REL-PAT.H09, REL-PAT.H10, REL-PAT.H11, REL-PAT.H12, REL-PAT.H13, REL-PAT.H14, REL-PAT.H15,
%   OBS: REL-PAT.N01, REL-PAT.N02, REL-PAT.N03, REL-PAT.N04, REL-PAT.N05, REL-PAT.N06, REL-PAT.K01,
%   OBS: REL-RDV.I01, REL-RDV.I02, REL-RDV.I03, REL-RDV.I04, REL-RDV.I05, REL-RDV.I06, REL-RDV.I07,
%   OBS: REL-RDV.I08, REL-RDV.I09, REL-RDV.I10, REL-RDV.R01, REL-RDV.R02, REL-RDV.R03, REL-RDV.R04,
%   OBS: REL-RDV.R05, REL-RDV.R06, REL-RDV.R07, REL-RDV.R08, REL-RDV.R09, REL-RDV.R10, REL-RDV.R11,
%   OBS: REL-RDV.R12, REL-RDV.R13, REL-RDV.R14, REL-RDV.C01, REL-RDV.C02, REL-RDV.C03, REL-RDV.C04,
%   OBS: REL-RDV.C05, REL-RDV.C06, REL-RDV.C07, REL-RDV.C08, REL-RDV.L01, REL-RDV.L02, REL-RDV.L03
%   HYP: (aucun)
%   CHI: (aucun)
%   SER: (aucune)
%   SECTIONS: 8

\chapter{Le système d'information hospitalier}

\section{Hosix.NET, origine et architecture}

Le système d'information en place est Hosix.NET, édité par SIVSA Soluciones Informáticas
\cite{s06}. Sa conception remonte aux années 1980, et l'éditeur en revendique le déploiement dans
plus de trente hôpitaux répartis sur six pays \cite{s07}.

L'origine du produit se lit dans sa nomenclature avant de se lire dans sa documentation
commerciale. La raison sociale de l'éditeur est espagnole, et les modules portent des noms
espagnols que la documentation reprend sans les traduire : \emph{Consultas} pour la gestion des
agendas et des rendez-vous \cite{s08}, \emph{Facturación} pour l'émission des factures
\cite{s09}, auxquels s'ajoutent \emph{Médicos}, \emph{Urgencias} et \emph{admisión}. Un produit
déployé au Maroc et servi en français conserve donc, dans son organisation interne, la langue de
son concepteur.

L'architecture fonctionnelle se contracte en trois paquets — Core, Clinic et Extension Packs —
qui rassemblent vingt-six modules : douze au Core, six au Clinic, huit aux Extension Packs
\cite{s06}. La répartition n'est pas indifférente. Le Core porte les fonctions administratives du
dossier — admission, agendas, facturation —, le Clinic les fonctions médicales, et les Extension
Packs des compléments contractables séparément. Un établissement peut donc exploiter la gestion
administrative des patients sans exploiter la couche clinique, et c'est une configuration que
l'observation conduite ici ne permet ni de confirmer ni d'exclure : le périmètre installé n'est
pas connu, la documentation de l'éditeur décrivant ce qui est contractable et jamais ce qui est
contracté.

Le produit a fait l'objet d'une migration technique, de WebForms vers le patron
modèle-vue-contrôleur, et l'éditeur la déclare achevée sur quatre modules : prescription,
admission, facturation et approvisionnements \cite{s07}. Cette migration n'est pas un détail
d'implantation. Elle explique qu'un même logiciel présente des écrans dont l'ergonomie diffère
d'un module à l'autre, et elle situe l'observation conduite ici dans un produit en cours de
transformation, non dans un état stable.

Une précision sur ce que ce chapitre ne dit pas, et elle vaut pour tout le dossier
documentaire. La documentation consultée est celle de l'éditeur, et elle est commerciale : elle
publie une fiche par module, et ces fiches énoncent des fonctions, pas des schémas de données.
Elles établissent qu'une fonction existe et ce qu'elle recouvre ; elles ne disent jamais quels
champs la portent, sous quel type, ni sous quelle contrainte. Aucune colonne du modèle construit
plus loin ne peut donc être justifiée par la seule documentation de l'éditeur.

Tout ce qui touche aux champs vient du relevé conduit au poste. C'est une matière d'une autre
nature : elle est datée, limitée à un profil et à quatre écrans, et elle ne couvre qu'une fraction
du produit. Elle a en revanche une propriété que la documentation n'a pas — elle porte sur
l'installation réelle de l'établissement, et non sur le produit générique. Les quatre sections de
relevé qui suivent la restituent intégralement, élément par élément.

\section{Les cinq profils applicatifs et les processus qu'ils portent}

L'écran de démarrage propose cinq profils applicatifs, et le relevé les reproduit sous leur
libellé exact : \texttt{MSM - FACTURATION SAA}, \texttt{MSM - GESTION DE RDV},
\texttt{MSM - RECOUVREMENT}, \texttt{MSM - URGENCE} et
\texttt{MSM RAPPORTS ET STATISTIQUES}\releve{REL-DEM.P02}. L'irrégularité du dernier — l'absence
de tiret séparateur — est celle de l'écran, et elle est conservée : un libellé se reproduit, il ne
se corrige pas.

Ces cinq profils ne sont pas cinq vues d'un même écran : chacun ouvre un jeu de modules distinct,
et les habilitations cloisonnent l'accès. Un seul a été ouvert au poste, les quatre autres n'ont
été relevés qu'à l'écran de démarrage. Le tableau ci-dessous rattache chaque profil au processus
qu'il porte et au module de l'éditeur qui le sert, en distinguant ce qui a été vu de ce qui a été
reconstruit par voie documentaire.

\begin{center}
\begin{tabular}{@{}p{4.3cm}p{4.6cm}p{3.1cm}p{2.2cm}@{}}
\toprule
Profil & Processus porté & Module & Établi par \\
\midrule
\texttt{MSM - GESTION DE RDV} & prise, modification et annulation de rendez-vous ; liste d'attente & \emph{Consultas} \cite{s08} & observation \\
\addlinespace
\texttt{MSM - FACTURATION SAA} & captation des mouvements, émission des factures & \emph{Facturación} \cite{s09} & documentation \\
\addlinespace
\texttt{MSM - RECOUVREMENT} & suivi des créances et des relances & \emph{Facturación} \cite{s09} & documentation \\
\addlinespace
\texttt{MSM - URGENCE} & admission et prise en charge aux urgences & \emph{Urgencias} & documentation \\
\addlinespace
\texttt{MSM RAPPORTS ET STATISTIQUES} & production des indicateurs remontés à la statistique nationale & — \cite{s10} & documentation \\
\bottomrule
\end{tabular}
\end{center}

\medskip

La documentation du module \emph{Consultas} énonce ce que le profil de rendez-vous recouvre :
génération d'agendas, prise de rendez-vous manuelle ou automatique, modification, annulation,
intégration à la liste d'attente, production de notifications, de listings et de statistiques
\cite{s08}. Celle de \emph{Facturación} énonce la captation manuelle ou automatique des
mouvements, le contrôle de l'état facturé ou en attente, le paramétrage par groupes, par types de
contrats et de prestations, et la facturation aux entités publiques, privées et mutuelles
\cite{s09}. L'éditeur ne publie aucune fiche distincte pour le recouvrement : la fonction est
intégrée au module de facturation, ce qui explique que deux profils s'y rattachent.

Le cinquième profil est le seul dont le contenu se déduise d'une source extérieure à l'éditeur :
la nomenclature des indicateurs qu'un hôpital nommé remonte à la statistique nationale
\cite{s10}.

\section{Méthode de relevé}

L'observation a porté sur un poste de travail du service, sur le profil
\texttt{MSM - GESTION DE RDV}, le 28 juillet 2026\releve{REL-DEM.P02}. Quatre écrans ont été
parcourus : l'écran de démarrage, la recherche d'identifiant patient, la fiche patient et l'écran
de prise de rendez-vous. Cent quarante-trois éléments y ont été relevés.

Chaque élément est consigné avec un identifiant stable, le bloc de l'écran auquel il appartient,
son libellé reproduit à l'identique, son type apparent, et la valeur affichée lorsqu'une valeur
l'était. L'identifiant n'est pas une commodité de rédaction : c'est la clé par laquelle le
registre de provenance des colonnes désigne son témoin d'observation, et par laquelle les
tableaux qui suivent se relient au modèle de données construit plus loin.

Deux conventions gouvernent la restitution, et toutes deux consistent à ne pas affirmer plus que
ce qui a été vu. La première : le caractère obligatoire ou facultatif d'un champ n'a été relevé
sur aucun écran, et n'est donc affirmé nulle part — une valeur binaire imposée ici produirait cent
quarante-trois affirmations non observées. La seconde : la mention \texttt{sans objet} désigne un
élément qui n'est pas un champ de saisie de dossier — bouton, onglet, entrée de menu, libellé de
profil, colonne de résultat, filtre de recherche. Les boutons des barres d'outils ne figurent pas
aux tableaux : ils déclenchent une opération et ne portent aucune donnée, et le relevé les déclare
comme tels.

\section{Relevé — l'écran de démarrage et les profils}

L'écran d'entrée identifie l'établissement et l'autorité dont il relève, puis propose les cinq
profils et quatre menus principaux. Onze éléments y sont relevés.

\begin{longtable}{@{}p{2.9cm}p{3.9cm}p{2.8cm}p{4.0cm}@{}}
\toprule
Identifiant & Libellé à l'écran & Type apparent & Valeur observée \\
\midrule\endfirsthead
\toprule
Identifiant & Libellé à l'écran & Type apparent & Valeur observée \\
\midrule\endhead
\multicolumn{4}{@{}l}{\textbf{Identification du site}} \\
\texttt{REL-DEM.I01} & Identification du site & texte & \texttt{005 Fès-Meknès; Sidi Said} \\
\texttt{REL-DEM.I02} & Ministère de tutelle & texte & \texttt{Ministère de la Santé} \\
\addlinespace
\multicolumn{4}{@{}l}{\textbf{Profils applicatifs}} \\
\texttt{REL-DEM.P01} & MSM - FACTURATION SAA & liste & — \\
\texttt{REL-DEM.P02} & MSM - GESTION DE RDV & liste & — \\
\texttt{REL-DEM.P03} & MSM - RECOUVREMENT & liste & — \\
\texttt{REL-DEM.P04} & MSM - URGENCE & liste & — \\
\texttt{REL-DEM.P05} & MSM RAPPORTS ET STATISTIQUES & liste & — \\
\addlinespace
\multicolumn{4}{@{}l}{\textbf{Menus principaux}} \\
\texttt{REL-DEM.M01} & Hosix.NET & menu & — \\
\texttt{REL-DEM.M02} & Données Patient & menu & — \\
\texttt{REL-DEM.M03} & Centre de Consultation & menu & — \\
\texttt{REL-DEM.M04} & Factures & menu & — \\
\addlinespace
\bottomrule
\end{longtable}

\medskip

Le code de région y est accolé au nom de l'établissement sous une graphie que le relevé
conserve telle quelle, séparateur compris\releve{REL-DEM.I01}. Les quatre menus donnent l'ossature
de la navigation, et deux d'entre eux mènent aux écrans relevés ci-après.

\section{Relevé — la recherche d'identifiant patient}

Cet écran est celui par lequel un agent retrouve un dossier avant toute opération. Trente-trois
éléments y sont relevés, répartis en quatre blocs.

\begin{longtable}{@{}p{2.9cm}p{3.9cm}p{2.8cm}p{4.0cm}@{}}
\toprule
Identifiant & Libellé à l'écran & Type apparent & Valeur observée \\
\midrule\endfirsthead
\toprule
Identifiant & Libellé à l'écran & Type apparent & Valeur observée \\
\midrule\endhead
\multicolumn{4}{@{}l}{\textbf{Onglets}} \\
\texttt{REL-IPP.O01} & Chercher & onglet & — \\
\texttt{REL-IPP.O02} & Chercher MPI & onglet & — \\
\texttt{REL-IPP.O03} & Chercher MPI Pondéré & onglet & — \\
\texttt{REL-IPP.O04} & Nouvel IPP & onglet & — \\
\addlinespace
\multicolumn{4}{@{}l}{\textbf{Critères de recherche}} \\
\texttt{REL-IPP.C01} & N° document & texte & \texttt{C.I.N} \\
\texttt{REL-IPP.C02} & Nom & texte & — \\
\texttt{REL-IPP.C03} & Prénom & texte & — \\
\texttt{REL-IPP.C04} & Date de naissance & date & — \\
\texttt{REL-IPP.C05} & Mois & entier & — \\
\texttt{REL-IPP.C06} & Année & entier & — \\
\texttt{REL-IPP.C07} & Approximation en années & entier & — \\
\texttt{REL-IPP.C08} & Âge & entier & — \\
\texttt{REL-IPP.C09} & Téléphone & texte & — \\
\texttt{REL-IPP.C10} & N° document de contact & texte & — \\
\texttt{REL-IPP.C11} & Sexe & liste & — \\
\texttt{REL-IPP.C12} & État civil & liste & — \\
\texttt{REL-IPP.C13} & N° assuré & texte & — \\
\addlinespace
\multicolumn{4}{@{}l}{\textbf{Filtres de population}} \\
\texttt{REL-IPP.F01} & Patients Hospitalisés & option & — \\
\texttt{REL-IPP.F02} & Patients aux Urgences & option & — \\
\texttt{REL-IPP.F03} & Tous les patients & option & — \\
\texttt{REL-IPP.F04} & Exitus & case\ a\ cocher & — \\
\addlinespace
\multicolumn{4}{@{}l}{\textbf{Colonnes de résultat}} \\
\texttt{REL-IPP.R01} & IPP & colonne & — \\
\texttt{REL-IPP.R02} & Nom & colonne & — \\
\texttt{REL-IPP.R03} & D. Naissance & colonne & — \\
\texttt{REL-IPP.R04} & N° pièce d'identité & colonne & — \\
\texttt{REL-IPP.R05} & N° Assu & colonne & — \\
\texttt{REL-IPP.R06} & E. Civil & colonne & — \\
\texttt{REL-IPP.R07} & Tél & colonne & — \\
\texttt{REL-IPP.R08} & Sexe & colonne & — \\
\texttt{REL-IPP.R09} & Province & colonne & — \\
\texttt{REL-IPP.R10} & Ville & colonne & — \\
\texttt{REL-IPP.R11} & Code postal & colonne & — \\
\texttt{REL-IPP.R12} & Probabilité & colonne & — \\
\addlinespace
\bottomrule
\end{longtable}

\medskip

\textbf{Un index maître des patients existe dans le système.} Le bloc d'onglets en porte la
preuve : à côté de la recherche ordinaire et de la création d'un nouvel identifiant, deux onglets
s'adressent explicitement à un index maître, dont l'un en recherche
pondérée\releve{REL-IPP.O03}. Le bloc des colonnes de résultat le confirme : la dernière colonne
affiche une \emph{probabilité}\releve{REL-IPP.R12}. Une recherche qui rend des probabilités n'est
pas une recherche exacte — c'est un rapprochement.

Ce que le relevé établit s'arrête là. Il atteste qu'un index maître et une recherche pondérée
existent dans l'installation observée ; il n'établit rien sur leur déploiement au-delà de cet
établissement, et aucune source du dossier ne le porte. La portée de cette observation pour la
modélisation est reprise au chapitre~\ref{chap:qualite}.

Les treize critères de recherche méritent une lecture. Ils ne cherchent pas seulement par
identifiant : nom, prénom, date de naissance, mais aussi mois et année seuls, âge, et une
\emph{approximation en années}. Un écran qui offre de chercher à l'âge approché est un écran conçu
pour des dossiers dont la date de naissance est incertaine ou absente.

\section{Relevé — la fiche patient}

La fiche patient est l'écran le plus dense du périmètre observé : cinquante-sept éléments, dont
dix boutons de barre d'outils qui ne figurent pas au tableau. Les quarante-sept restants se
répartissent en trois onglets et cinq blocs de données.

\begin{center}
\begin{tabular}{@{}c@{}}
\fbox{\begin{minipage}{0.86\linewidth}
\centering\small
\fbox{\begin{minipage}{0.95\linewidth}\centering Barre d'actions \\ \footnotesize dix boutons, hors tableau\end{minipage}}\\[3pt]
\fbox{\begin{minipage}{0.95\linewidth}\centering Onglets : Général \quad|\quad Contacts \quad|\quad Quota soins primaires\end{minipage}}\\[3pt]
\fbox{\begin{minipage}{0.95\linewidth}\centering Données principales \\ \footnotesize seize éléments\end{minipage}}\\[3pt]
\fbox{\begin{minipage}{0.46\linewidth}\centering Compagnie d'assurance \\ \footnotesize six éléments\end{minipage}}%
\hfill
\fbox{\begin{minipage}{0.46\linewidth}\centering Né \\ \footnotesize six éléments\end{minipage}}\\[3pt]
\fbox{\begin{minipage}{0.95\linewidth}\centering Domicile \\ \footnotesize quinze éléments\end{minipage}}\\[3pt]
\fbox{\begin{minipage}{0.95\linewidth}\centering Commentaire\end{minipage}}
\end{minipage}}
\end{tabular}
\end{center}

\begin{center}
\footnotesize Agencement des blocs de la fiche patient. Reconstitution structurelle : aucun
contenu, aucune valeur.
\end{center}

\medskip

\begin{longtable}{@{}p{2.9cm}p{3.9cm}p{2.8cm}p{4.0cm}@{}}
\toprule
Identifiant & Libellé à l'écran & Type apparent & Valeur observée \\
\midrule\endfirsthead
\toprule
Identifiant & Libellé à l'écran & Type apparent & Valeur observée \\
\midrule\endhead
\multicolumn{4}{@{}l}{\textbf{Onglets}} \\
\texttt{REL-PAT.T01} & Général & onglet & — \\
\texttt{REL-PAT.T02} & Contacts & onglet & — \\
\texttt{REL-PAT.T03} & Quota soins primaires & onglet & — \\
\addlinespace
\multicolumn{4}{@{}l}{\textbf{Données principales}} \\
\texttt{REL-PAT.D01} & N° IPP & texte & — \\
\texttt{REL-PAT.D02} & Nom & texte & — \\
\texttt{REL-PAT.D03} & Nom de famille 1 & texte & — \\
\texttt{REL-PAT.D04} & Nom de famille 2 & texte & — \\
\texttt{REL-PAT.D05} & Sexe & liste & — \\
\texttt{REL-PAT.D06} & D. Nai & date & — \\
\texttt{REL-PAT.D07} & Âge & entier & — \\
\texttt{REL-PAT.D08} & Num & non\ determine & — \\
\texttt{REL-PAT.D09} & Type pièce d'identité & liste & \texttt{C} / \texttt{C.I.N} \\
\texttt{REL-PAT.D10} & N° pièce d'identité & texte & — \\
\texttt{REL-PAT.D11} & E. Civil & liste & \texttt{C} / \texttt{CÉLIBATAIRE} \\
\texttt{REL-PAT.D12} & Type patient & liste & — \\
\texttt{REL-PAT.D13} & Date photo & date & — \\
\texttt{REL-PAT.D14} & Modifié par & texte & — \\
\texttt{REL-PAT.D15} & Créé par & texte & — \\
\texttt{REL-PAT.D16} & Date d'attribution & date & — \\
\addlinespace
\multicolumn{4}{@{}l}{\textbf{Compagnie d'assurance}} \\
\texttt{REL-PAT.A01} & Compagnie d'assur. & liste & \texttt{00116} / \texttt{CNSS PIPC} \\
\texttt{REL-PAT.A02} & Police & texte & — \\
\texttt{REL-PAT.A03} & N° Assu & texte & — \\
\texttt{REL-PAT.A04} & Profession & texte & — \\
\texttt{REL-PAT.A05} & Num. inscription & texte & — \\
\texttt{REL-PAT.A06} & Date inscription & date & — \\
\addlinespace
\multicolumn{4}{@{}l}{\textbf{Domicile}} \\
\texttt{REL-PAT.H01} & Type & liste & — \\
\texttt{REL-PAT.H02} & Adresse & texte & — \\
\texttt{REL-PAT.H03} & Code postal & texte & — \\
\texttt{REL-PAT.H04} & État & liste & — \\
\texttt{REL-PAT.H05} & Ville & liste & — \\
\texttt{REL-PAT.H06} & Quartier & texte & — \\
\texttt{REL-PAT.H07} & Nationalité & liste & \texttt{504} / \texttt{MAROC} \\
\texttt{REL-PAT.H08} & Téléphone 1 & texte & — \\
\texttt{REL-PAT.H09} & Téléphone 2 & texte & — \\
\texttt{REL-PAT.H10} & Téléphone 3 & texte & — \\
\texttt{REL-PAT.H11} & Téléphone 4 & texte & — \\
\texttt{REL-PAT.H12} & Avertissements SMS & case\ a\ cocher & — \\
\texttt{REL-PAT.H13} & E-mail & texte & — \\
\texttt{REL-PAT.H14} & Avertissements e-mail & case\ a\ cocher & — \\
\texttt{REL-PAT.H15} & Environnement & liste & — \\
\addlinespace
\multicolumn{4}{@{}l}{\textbf{Né}} \\
\texttt{REL-PAT.N01} & Nom. Père & texte & — \\
\texttt{REL-PAT.N02} & Nom. Mère & texte & — \\
\texttt{REL-PAT.N03} & Lieu de naissance - État & liste & — \\
\texttt{REL-PAT.N04} & Ville & liste & — \\
\texttt{REL-PAT.N05} & Pays & liste & \texttt{504} / \texttt{MAROC} \\
\texttt{REL-PAT.N06} & Quartier & texte & — \\
\addlinespace
\multicolumn{4}{@{}l}{\textbf{Commentaire}} \\
\texttt{REL-PAT.K01} & Commentaire & zone\ texte & — \\
\addlinespace
\bottomrule
\end{longtable}

\medskip

Quatre valeurs affichées documentent la forme des nomenclatures du
système\releve{REL-PAT.D09}. Deux sont des couples code-libellé : la compagnie d'assurance et la
nationalité, cette dernière relevée deux fois — au domicile et au lieu de naissance — sous le même
couple. Deux autres sont des couples de valeurs codées : le type de pièce d'identité et l'état
civil, tous deux abrégés par une lettre suivie de son libellé développé. Le système code donc ses
références et affiche les deux faces du couple ; c'est ce qui permet, plus loin, de reconstituer
des dimensions sans inventer de libellés.

Le bloc du domicile porte quatre lignes de téléphone et deux cases d'avertissement, l'une par
message court, l'autre par courrier électronique\releve{REL-PAT.H12}. Les canaux de contact
existent donc déjà dans la fiche.

\section{Relevé — donner rendez-vous}

L'écran de prise de rendez-vous compte quarante-deux éléments, dont sept boutons de barre d'outils
hors tableau. Les trente-cinq restants se répartissent en quatre blocs : un rappel d'identité du
patient, le rendez-vous lui-même, un bloc de contrôle des modifications, et la liste d'attente des
consultations.

\begin{center}
\begin{tabular}{@{}c@{}}
\fbox{\begin{minipage}{0.86\linewidth}
\centering\small
\fbox{\begin{minipage}{0.95\linewidth}\centering Barre d'actions \\ \footnotesize sept boutons, hors tableau\end{minipage}}\\[3pt]
\fbox{\begin{minipage}{0.95\linewidth}\centering Identification du patient \\ \footnotesize dix éléments, rappelés de la fiche\end{minipage}}\\[3pt]
\fbox{\begin{minipage}{0.95\linewidth}\centering Rendez-vous \\ \footnotesize quatorze éléments\end{minipage}}\\[3pt]
\fbox{\begin{minipage}{0.46\linewidth}\centering Contrôle de modifications \\ \footnotesize huit éléments\end{minipage}}%
\hfill
\fbox{\begin{minipage}{0.46\linewidth}\centering Liste d'attente \\ \footnotesize trois colonnes\end{minipage}}
\end{minipage}}
\end{tabular}
\end{center}

\begin{center}
\footnotesize Agencement des blocs de l'écran de prise de rendez-vous. Reconstitution
structurelle : aucun contenu, aucune valeur.
\end{center}

\medskip

\begin{longtable}{@{}p{2.9cm}p{3.9cm}p{2.8cm}p{4.0cm}@{}}
\toprule
Identifiant & Libellé à l'écran & Type apparent & Valeur observée \\
\midrule\endfirsthead
\toprule
Identifiant & Libellé à l'écran & Type apparent & Valeur observée \\
\midrule\endhead
\multicolumn{4}{@{}l}{\textbf{Identification du patient}} \\
\texttt{REL-RDV.I01} & N° IPP & texte & — \\
\texttt{REL-RDV.I02} & Téléphone & texte & — \\
\texttt{REL-RDV.I03} & Compagnie & liste & — \\
\texttt{REL-RDV.I04} & N° Assu & texte & — \\
\texttt{REL-RDV.I05} & Âge & entier & — \\
\texttt{REL-RDV.I06} & Date naissance & date & — \\
\texttt{REL-RDV.I07} & C.I.N & texte & — \\
\texttt{REL-RDV.I08} & Num & non\ determine & — \\
\texttt{REL-RDV.I09} & Adresse & texte & — \\
\texttt{REL-RDV.I10} & Ville & liste & — \\
\addlinespace
\multicolumn{4}{@{}l}{\textbf{Rendez-vous}} \\
\texttt{REL-RDV.R01} & Agenda & liste & — \\
\texttt{REL-RDV.R02} & Activité & liste & — \\
\texttt{REL-RDV.R03} & Origine & liste & \texttt{AU} / \texttt{AUTRES HÔPITAUX} \\
\texttt{REL-RDV.R04} & Hôpital/C.S. & liste & — \\
\texttt{REL-RDV.R05} & Médecin ext. & texte & — \\
\texttt{REL-RDV.R06} & Service ext. & texte & — \\
\texttt{REL-RDV.R07} & Observations & zone\ texte & — \\
\texttt{REL-RDV.R08} & Date rendez-vous & horodatage & \texttt{07/28/2026 10:57:07 AM} \\
\texttt{REL-RDV.R09} & Rendez-vous supplémentaire & case\ a\ cocher & — \\
\texttt{REL-RDV.R10} & Type d'attention & liste & — \\
\texttt{REL-RDV.R11} & État & liste & \texttt{En instance} \\
\texttt{REL-RDV.R12} & Durée & entier & \texttt{0} \\
\texttt{REL-RDV.R13} & Date réception & horodatage & — \\
\texttt{REL-RDV.R14} & Imprimer données & case\ a\ cocher & — \\
\addlinespace
\multicolumn{4}{@{}l}{\textbf{Contrôle de modifications}} \\
\texttt{REL-RDV.C01} & Créé par & texte & — \\
\texttt{REL-RDV.C02} & Date création & horodatage & \texttt{07/28/2026 10:57:06 AM} \\
\texttt{REL-RDV.C03} & Modifié par & texte & — \\
\texttt{REL-RDV.C04} & Date mod. & horodatage & — \\
\texttt{REL-RDV.C05} & Confirmé par & texte & — \\
\texttt{REL-RDV.C06} & Date conf. & horodatage & — \\
\texttt{REL-RDV.C07} & Annulé par & texte & — \\
\texttt{REL-RDV.C08} & Date annul. & horodatage & — \\
\addlinespace
\multicolumn{4}{@{}l}{\textbf{Liste d'attente des consultations}} \\
\texttt{REL-RDV.L01} & Service & colonne & — \\
\texttt{REL-RDV.L02} & Agenda & colonne & — \\
\texttt{REL-RDV.L03} & Activité & colonne & — \\
\addlinespace
\bottomrule
\end{longtable}

\medskip

Le bloc du rendez-vous porte l'origine de la demande sous forme de couple code-libellé, et la
valeur observée désigne un adressage par une autre structure
hospitalière\releve{REL-RDV.R03}. L'état du rendez-vous à la création et sa durée sont également
lisibles\releve{REL-RDV.R11}. Le bloc de contrôle des modifications porte quatre couples
d'auteur et d'horodatage — création, modification, confirmation, annulation — dont la portée est
traitée à la section suivante.

Les dix éléments du bloc d'identification ne sont pas des champs de saisie propres à cet écran :
ils rappellent des valeurs de la fiche patient. Le modèle de données ne les enregistre donc pas
deux fois.

\section{Trois observations exploitables}

Trois relevés de cet écran dépassent la description : ils fondent une grandeur, une distinction et
un diagnostic que le reste du rapport reprend. Chacun est une observation, et rien d'autre.

\textbf{L'écart entre la création et le rendez-vous est le délai
d'obtention}\releve{REL-RDV.C02}. Le bloc de contrôle des modifications horodate la création ; le
bloc du rendez-vous horodate la date convenue. Leur différence est le délai qu'un patient attend,
grandeur que ni l'un ni l'autre des deux champs ne porte seul. Sur le rendez-vous observé, cette
différence vaut une seconde : la création à \texttt{10:57:06} et le rendez-vous à
\texttt{10:57:07} du même jour, soit un rendez-vous donné pour l'instant même. Le cas est extrême
et il est instructif : il montre que le système accepte un délai nul, donc que la grandeur est
bornée à zéro et non strictement positive. Le chapitre~\ref{chap:activite} en fait un indicateur.

\textbf{Un couple de champs sépare l'annulation de l'absence}\releve{REL-RDV.C07}. Le bloc de
contrôle porte un champ d'auteur de confirmation et un champ d'auteur d'annulation, chacun avec
son horodatage. Sans ce couple, un rendez-vous non honoré serait un état unique. Avec lui, deux
situations se distinguent : l'annulation déclarée, qui laisse au service le temps de réaffecter le
créneau, et l'absence non prévenue, qui le perd. La distinction n'est pas cosmétique — elle sépare
deux taux dont les remèdes diffèrent, et le chapitre~\ref{chap:activite} les mesure séparément.

\textbf{La recherche pondérée est un aveu}\releve{REL-IPP.O03}. Un système qui offre de chercher
un patient par similarité, et qui affiche une probabilité en regard de chaque résultat, admet que
la recherche exacte ne suffit pas. Elle ne suffit pas parce qu'une même personne y possède
plusieurs dossiers. La fonction existe dans le produit installé ; ce que le relevé ne dit pas,
c'est si le service l'emploie. Le chapitre~\ref{chap:qualite} construit un rapprochement
probabiliste sur ce constat, et le confronte à une collision exacte pour mesurer ce que la
similarité apporte.

Une quatrième observation, mineure et têtue, court à travers les trois précédentes : les
horodatages s'affichent au format mois-jour-année avec un indicateur de méridien, dans une
interface par ailleurs française\releve{REL-RDV.R08}. La date observée le montre sans ambiguïté —
le quantième vaut vingt-huit, ce qui exclut toute lecture jour-mois, et lève l'incertitude qu'une
date antérieure au treizième jour du mois aurait laissée entière.

Ce détail d'affichage n'est pas cosmétique. Une extraction du système rend ses dates dans la forme
où l'écran les compose ; toute date entrant dans la chaîne de données doit donc être convertie
depuis ce format, et la conversion doit être écrite une fois et vérifiée, faute de quoi un jour et
un mois s'échangent silencieusement onze fois sur douze sans qu'aucune valeur ne devienne
invalide. C'est une contrainte de chargement, reprise au chapitre~\ref{chap:architecture}.

Ces quatre observations partagent une propriété qu'il faut nommer, parce qu'elle gouverne la
suite. Aucune ne résulte d'un entretien ni d'une documentation : chacune est lisible sur un écran,
et chacune est consignée sous un identifiant que le lecteur peut retrouver dans les tableaux des
sections précédentes. Elles ne disent rien de la fréquence des situations qu'elles révèlent — un
seul rendez-vous a été observé, et un cas n'est pas une distribution. Ce qu'elles établissent est
qu'une grandeur est calculable, qu'une distinction est enregistrée, qu'une fonction existe et
qu'un format s'impose. C'est précisément ce dont la modélisation a besoin : non pas des valeurs,
mais la certitude que les champs qui les porteraient existent.
